% =====================================================================
%  Banco complementar --- Circuitos Paralelo  (REVISADO)
%  Substitui a versao gerada por template (11 clones em N2, 13 em N3,
%  9 questoes conceituais indevidamente marcadas como N4).
%  Sem sobreposicao com o cap03 (n resistores iguais, R desconhecido a
%  partir de Req, divisor 12 A entre 4 e 12 ohm, tres lampadas de 60 W,
%  fio ideal em paralelo, banco de 100 ohm/2 W e a demonstracao de que
%  Req < menor resistor ja estao la).
%  Distribuicao mantida: 7 N1 + 11 N2 + 13 N3 + 9 N4 = 40
% =====================================================================

\subsubsection*{Banco complementar --- Circuitos paralelo}

\begin{atencao}[Organização do banco]
No paralelo, a \textbf{tensão é comum} e as \textbf{correntes se somam}. Daí
decorre que corrente e potência se distribuem \emph{inversamente} a $R$ ---
exatamente o oposto da série. Trabalhe com \textbf{condutâncias}: elas somam
diretamente e tornam quase toda conta trivial. O N3 explora falhas de ramo,
fuga, deriva térmica e instrumentação; o N4 trata de projeto de \emph{shunt},
sensibilidade, propagação de incerteza e seletividade de proteção.
\end{atencao}

\blocofixacao

\begin{questao}[1]{}
Determine a resistência equivalente de $R_1=\SI{30}{\ohm}$ e
$R_2=\SI{60}{\ohm}$ em paralelo.
\resposta{$R_{eq}=\SI{20}{\ohm}$}
\resolucao{Para dois resistores, a regra do produto pela soma é a mais rápida:
\[
R_{eq}=\frac{R_1R_2}{R_1+R_2}=\frac{30\cdot60}{90}=\SI{20}{\ohm}.
\]
\textbf{Conferência instantânea:} o resultado ($\SI{20}{\ohm}$) é menor que o
menor dos dois ($\SI{30}{\ohm}$) --- se não for, há erro. Cuidado: a regra
produto/soma vale \emph{apenas} para dois resistores.}
\end{questao}

\begin{questao}[1]{}
Três resistores de $\SI{4}{\ohm}$, $\SI{8}{\ohm}$ e $\SI{8}{\ohm}$ estão em
paralelo. Calcule a condutância total e $R_{eq}$.
\resposta{$G=\SI{0.5}{\siemens}$; $R_{eq}=\SI{2}{\ohm}$}
\resolucao{
\[
G=\frac14+\frac18+\frac18=0{,}25+0{,}125+0{,}125=\SI{0.5}{\siemens}
\;\Longrightarrow\;
R_{eq}=\frac{1}{0{,}5}=\SI{2}{\ohm}.
\]
Com três ou mais ramos, some as condutâncias e inverta \emph{uma única vez} no
final. Inverter parcialmente a cada passo é a origem mais comum de erro.}
\end{questao}

\begin{questao}[1]{}
Três resistores iguais de $\SI{90}{\ohm}$ são associados em paralelo. Determine
$R_{eq}$.
\resposta{$R_{eq}=\SI{30}{\ohm}$}
\resolucao{Para $n$ resistores \emph{iguais},
\[
R_{eq}=\frac{R}{n}=\frac{90}{3}=\SI{30}{\ohm}.
\]
O atalho vale só quando todos são iguais; com valores diferentes é obrigatório
somar condutâncias.}
\end{questao}

\begin{questao}[1]{}
Acrescentando-se mais um ramo \textbf{em paralelo} a um circuito alimentado por
fonte ideal de tensão, a corrente total:
\begin{alternativas}[2]
\item aumenta
\item diminui
\item não se altera
\item pode aumentar ou diminuir, conforme o valor
\end{alternativas}
\resposta{(a)}
\resolucao{Cada novo ramo acrescenta uma condutância positiva:
$G_{eq}=\sum G_k$ só pode crescer, logo $R_{eq}$ só pode diminuir. Com $V$
fixo, $I=V/R_{eq}$ aumenta.\\
Note o contraste com a série, em que acrescentar elemento \emph{reduz} a
corrente. E note também que os ramos já existentes continuam com \emph{a mesma}
corrente --- quem cresce é o total.}
\end{questao}

\begin{questao}[1]{}
Um ramo de $\SI{48}{\ohm}$ pertence a uma associação paralela ligada a
$\SI{24}{\volt}$. Determine a corrente desse ramo.
\resposta{$I=\SI{0.5}{\ampere}$}
\resolucao{
\[
I=\frac{V}{R}=\frac{24}{48}=\SI{0.5}{\ampere}.
\]
No paralelo, a tensão é \emph{comum} a todos os ramos: a corrente de cada um
depende só do próprio resistor. Não é preciso conhecer os demais ramos nem
$R_{eq}$.}
\end{questao}

\begin{questao}[1]{}
Uma associação paralela sob $\SI{40}{\volt}$ tem ramos que conduzem
$\SI{2}{\ampere}$, $\SI{3}{\ampere}$ e $\SI{5}{\ampere}$. Determine a corrente
total e $R_{eq}$.
\resposta{$I_t=\SI{10}{\ampere}$; $R_{eq}=\SI{4}{\ohm}$}
\resolucao{Pela LKC, as correntes de ramo se somam:
$I_t=2+3+5=\SI{10}{\ampere}$.
\[
R_{eq}=\frac{V}{I_t}=\frac{40}{10}=\SI{4}{\ohm}.
\]
Observe que não foi preciso conhecer os resistores individualmente --- embora
eles saiam de imediato: $20$, $13{,}3$ e $\SI{8}{\ohm}$.}
\end{questao}

\begin{questao}[1]{}
Em uma associação paralela, o ramo que conduz a \textbf{maior} corrente é:
\begin{alternativas}[2]
\item o de maior resistência
\item o de menor resistência
\item o mais próximo da fonte
\item todos conduzem igualmente
\end{alternativas}
\resposta{(b)}
\resolucao{Com $V$ comum, $I_k=V/R_k$: a corrente é \emph{inversamente}
proporcional à resistência. O ramo de menor $R$ é o caminho de menor oposição e
leva a maior corrente.\\
A alternativa (c) é sedutora, mas falsa: em circuitos de parâmetros
concentrados, a posição física dos ramos é irrelevante --- só importa a
topologia.}
\end{questao}

\blocoaplicacao

\begin{questao}[2]{}
$R_1=\SI{10}{\ohm}$, $R_2=\SI{20}{\ohm}$ e $R_3=\SI{20}{\ohm}$ estão em
paralelo sob $\SI{60}{\volt}$. Determine as correntes de ramo, a corrente total
e $R_{eq}$.
\resposta{$\SI{6}{\ampere}$, $\SI{3}{\ampere}$, $\SI{3}{\ampere}$;
$I_t=\SI{12}{\ampere}$; $R_{eq}=\SI{5}{\ohm}$}
\resolucao{Tensão comum, então cada ramo é independente:
\[
I_1=\frac{60}{10}=6,\quad I_2=I_3=\frac{60}{20}=\SI{3}{\ampere}.
\]
\[
I_t=6+3+3=\SI{12}{\ampere}
\;\Longrightarrow\;
R_{eq}=\frac{60}{12}=\SI{5}{\ohm}.
\]
\textbf{Conferência:} $G=0{,}1+0{,}05+0{,}05=\SI{0.2}{\siemens}$ e
$1/0{,}2=\SI{5}{\ohm}$ \checkmark.}
\end{questao}

\begin{questao}[2]{}
No circuito da questão anterior, calcule a potência de cada ramo e confronte
com a potência total.
\resposta{$\SI{360}{\watt}$, $\SI{180}{\watt}$, $\SI{180}{\watt}$; total
$\SI{720}{\watt}$}
\resolucao{Com $V$ comum, use $P=V^{2}/R$:
\[
P_1=\frac{3600}{10}=360,\qquad P_2=P_3=\frac{3600}{20}=\SI{180}{\watt}.
\]
Fonte: $P=VI_t=60\cdot12=\SI{720}{\watt}=360+180+180$ \checkmark.\\
\textbf{Regra do paralelo:} $P\propto1/R$ --- o \emph{menor} resistor dissipa
mais. Compare com a série, em que $P\propto R$. Trocar essas duas regras é o
erro conceitual mais frequente em associações.}
\end{questao}

\begin{questao}[2]{}
Uma corrente de $\SI{15}{\ampere}$ chega a um nó e se divide entre
$R_1=\SI{8}{\ohm}$ e $R_2=\SI{12}{\ohm}$. Determine as correntes de ramo pela
fórmula do divisor.
\resposta{$I_1=\SI{9}{\ampere}$, $I_2=\SI{6}{\ampere}$}
\resolucao{Para \emph{dois} ramos, a corrente de um deles é proporcional à
resistência do \textbf{outro}:
\[
I_1=I_t\frac{R_2}{R_1+R_2}=15\cdot\frac{12}{20}=\SI{9}{\ampere},
\qquad
I_2=15\cdot\frac{8}{20}=\SI{6}{\ampere}.
\]
\textbf{Verificação:} $9+6=\SI{15}{\ampere}$ \checkmark, e a maior corrente ficou
no \emph{menor} resistor \checkmark. O "cruzamento" dos índices é o ponto que
mais confunde --- ele existe porque $I_k\propto G_k$, e escrever a fórmula em
condutâncias ($I_1=I_tG_1/(G_1+G_2)$) elimina a inversão.}
\end{questao}

\begin{questao}[2]{}
Três resistores em paralelo sob $\SI{100}{\volt}$ drenam $\SI{5}{\ampere}$ no
total. Dois deles valem $\SI{50}{\ohm}$ e $\SI{100}{\ohm}$. Determine o
terceiro.
\resposta{$R_3=\SI{50}{\ohm}$}
\resolucao{Trabalhe em condutâncias, onde a subtração é direta:
\[
G_{eq}=\frac{I_t}{V}=\frac{5}{100}=\SI{0.05}{\siemens},
\qquad
G_1+G_2=0{,}02+0{,}01=\SI{0.03}{\siemens}.
\]
\[
G_3=0{,}05-0{,}03=\SI{0.02}{\siemens}
\;\Longrightarrow\;
R_3=\SI{50}{\ohm}.
\]
Tentar subtrair \emph{resistências} aqui levaria a um resultado sem sentido:
no paralelo, o que é aditivo é $G$, não $R$.}
\end{questao}

\begin{questao}[2]{}
Duas lâmpadas de resistências $\SI{240}{\ohm}$ e $\SI{960}{\ohm}$ são ligadas
\textbf{em paralelo} a $\SI{120}{\volt}$. Qual brilha mais? Compare com o
resultado da mesma dupla em série.
\resposta{Em paralelo, a de $\SI{240}{\ohm}$ ($\SI{60}{\watt}$ contra
$\SI{15}{\watt}$) --- inverte-se em relação à série}
\resolucao{
\[
P_{240}=\frac{120^{2}}{240}=\SI{60}{\watt},
\qquad
P_{960}=\frac{120^{2}}{960}=\SI{15}{\watt}.
\]
Em paralelo, $V$ é comum e $P\propto1/R$: brilha mais a de \emph{menor}
resistência. Em série (mesma dupla), $I$ é comum e $P\propto R$: brilharia mais
a de $\SI{960}{\ohm}$.\\
\textbf{Conclusão que vale guardar:} não existe "a lâmpada mais potente" em
termos absolutos --- a potência depende da \emph{ligação}. É por isso que
lâmpadas residenciais, sempre em paralelo, obedecem à intuição usual
($\SI{100}{\watt}$ brilha mais que $\SI{25}{\watt}$), mas guirlandas em série
se comportam ao contrário.}
\end{questao}

\begin{questao}[2]{}
Um galvanômetro tem resistência interna $\SI{100}{\ohm}$ e deflexão total com
$\SI{1}{\milli\ampere}$. Determine o resistor \emph{shunt} que amplia a escala
para $\SI{11}{\milli\ampere}$.
\resposta{$R_{sh}=\SI{10}{\ohm}$}
\resolucao{O shunt fica em paralelo com o galvanômetro e desvia o excedente:
$I_{sh}=11-1=\SI{10}{\milli\ampere}$. Como a tensão é comum,
\[
R_{sh}I_{sh}=R_gI_g
\;\Longrightarrow\;
R_{sh}=\frac{R_gI_g}{I_{sh}}=\frac{100\cdot1}{10}=\SI{10}{\ohm}.
\]
A regra geral é $R_{sh}=\dfrac{R_g}{n-1}$, com $n=I_{total}/I_g$ o fator de
ampliação --- aqui $n=11$ e $R_{sh}=100/10$ \checkmark.}
\end{questao}

\begin{questao}[2]{}
Um circuito residencial de $\SI{220}{\volt}$ alimenta, em paralelo, cargas de
$\SI{4400}{\watt}$, $\SI{1100}{\watt}$ e $\SI{550}{\watt}$. Determine a
corrente total e escolha o disjuntor.
\resposta{$I_t=\SI{27.5}{\ampere}$; disjuntor de $\SI{32}{\ampere}$}
\resolucao{Cada carga drena $I=P/V$:
\[
I_1=\frac{4400}{220}=20,\quad
I_2=\frac{1100}{220}=5,\quad
I_3=\frac{550}{220}=\SI{2.5}{\ampere}.
\]
Total: $\SI{27.5}{\ampere}$ (as correntes se somam porque as cargas estão em
paralelo).\\
O disjuntor comercial imediatamente superior é o de $\SI{32}{\ampere}$ ---
$\SI{25}{\ampere}$ seria insuficiente. \textbf{Atenção:} escolhido o disjuntor,
a seção do condutor deve suportar \emph{a corrente do disjuntor}, não a da
carga; caso contrário, o cabo vira o elo fraco da proteção.}
\end{questao}

\begin{questao}[2]{}
Dois condutores idênticos de $\SI{0.8}{\ohm}$ cada são ligados em paralelo para
alimentar uma carga. Determine a resistência resultante e explique a prática.
\resposta{$\SI{0.4}{\ohm}$ --- metade}
\resolucao{$R_{eq}=0{,}8/2=\SI{0.4}{\ohm}$.\\
Dobrar o número de condutores idênticos equivale a dobrar a seção: a
resistência cai à metade, e com ela a queda de tensão e as perdas
$RI^{2}$ na linha.\\
\textbf{Condição prática:} os condutores devem ter mesmo comprimento, seção e
material. Se um for mais curto, ele terá menor resistência e assumirá corrente
desproporcional --- podendo sobreaquecer enquanto o outro fica ocioso. Por isso
condutores em paralelo são sempre especificados com o mesmo lance.}
\end{questao}

\begin{questao}[2]{}
Um resistor de $\SI{1}{\mega\ohm}$ é ligado em paralelo com um de
$\SI{100}{\ohm}$. A resistência equivalente:
\begin{alternativas}[2]
\item fica praticamente igual a $\SI{100}{\ohm}$
\item fica praticamente igual a $\SI{1}{\mega\ohm}$
\item vale aproximadamente $\SI{500}{\kilo\ohm}$
\item vale exatamente $\SI{50}{\ohm}$
\end{alternativas}
\resposta{(a)}
\resolucao{$G=0{,}01+10^{-6}\approx\SI{0.010001}{\siemens}$, logo
$R_{eq}\approx\SI{99.99}{\ohm}$ --- diferença de $0{,}01\%$.\\
\textbf{Princípio geral:} no paralelo, quem manda é o \emph{menor} resistor;
ramos muito maiores são praticamente circuitos abertos. É exatamente por isso
que um voltímetro de alta impedância "não carrega" o circuito sob medição.}
\end{questao}

\begin{questao}[2]{}
Um chuveiro possui duas resistências de $\SI{22}{\ohm}$ que podem ser ligadas
isoladamente ou em paralelo, sob $\SI{220}{\volt}$. Determine a potência em
cada caso.
\resposta{$\SI{2200}{\watt}$ com uma; $\SI{4400}{\watt}$ com as duas}
\resolucao{Uma resistência:
$P=\dfrac{220^{2}}{22}=\SI{2200}{\watt}$.\\
Duas em paralelo: $R_{eq}=\SI{11}{\ohm}$ e
$P=\dfrac{220^{2}}{11}=\SI{4400}{\watt}$.\\
A potência \emph{dobra} porque a resistência caiu à metade sob a mesma tensão.
É assim que funcionam as chaves inverno/verão. \textbf{Consequência para a
instalação:} a corrente também dobra ($10\to\SI{20}{\ampere}$), e o
dimensionamento de disjuntor e condutor deve ser feito para a posição de maior
potência.}
\end{questao}

\begin{questao}[2]{}
Em uma associação paralela mediram-se as correntes de ramo $\SI{4}{\ampere}$,
$\SI{1}{\ampere}$ e $\SI{1}{\ampere}$, com corrente total de
$\SI{6}{\ampere}$ e tensão de $\SI{20}{\volt}$. Verifique a consistência e
determine as três resistências e a potência total.
\resposta{Consistente; $\SI{5}{\ohm}$, $\SI{20}{\ohm}$, $\SI{20}{\ohm}$;
$P=\SI{120}{\watt}$}
\resolucao{\textbf{LKC:} $4+1+1=\SI{6}{\ampere}$ \checkmark.\\
\textbf{Resistências:} $R_k=V/I_k$ $\Rightarrow$ $20/4=5$ e
$20/1=\SI{20}{\ohm}$ (duas vezes).\\
\textbf{Potência:} $P=VI_t=20\cdot6=\SI{120}{\watt}$; por ramo,
$80+20+20=\SI{120}{\watt}$ \checkmark.\\
Dois testes independentes --- LKC e balanço de potência --- fecham. Se apenas o
primeiro fechasse, o erro estaria na leitura da tensão.}
\end{questao}

\blocoaprofundamento

\begin{questao}[3]{}
Uma corrente total de $\SI{6}{\ampere}$ divide-se entre $\SI{5}{\ohm}$,
$\SI{20}{\ohm}$ e $\SI{20}{\ohm}$ em paralelo.
\begin{itensq}
\item Determine $R_{eq}$ e a tensão da associação.
\item Obtenha as três correntes pelo divisor em condutâncias.
\item Verifique por dois caminhos independentes.
\end{itensq}
\resposta{$R_{eq}=\SI{3.33}{\ohm}$, $V=\SI{20}{\volt}$;
$\SI{4}{\ampere}$, $\SI{1}{\ampere}$, $\SI{1}{\ampere}$}
\resolucao{(a) $G=0{,}2+0{,}05+0{,}05=\SI{0.3}{\siemens}$,
$R_{eq}=1/0{,}3=\SI{3.33}{\ohm}$ e $V=I_tR_{eq}=6\cdot3{,}33=\SI{20}{\volt}$.\\
(b) Com três ou mais ramos, a forma correta do divisor usa condutâncias:
\[
I_k=I_t\frac{G_k}{G_{eq}}
\;\Rightarrow\;
I_1=6\cdot\frac{0{,}2}{0{,}3}=4,\quad
I_2=I_3=6\cdot\frac{0{,}05}{0{,}3}=\SI{1}{\ampere}.
\]
\textbf{Alerta:} a fórmula $I_1=I_tR_2/(R_1+R_2)$ \emph{não} se generaliza para
três ramos --- ela é um caso particular que só funciona com dois.\\
(c) \emph{LKC:} $4+1+1=6$ \checkmark. \emph{Lei de Ohm por ramo:}
$20/5=4$ e $20/20=1$ \checkmark.}
\end{questao}

\begin{questao}[3]{}
Três ramos de $\SI{10}{\ohm}$, $\SI{20}{\ohm}$ e $\SI{20}{\ohm}$ operam sob
fonte \textbf{ideal} de $\SI{60}{\volt}$. O ramo de $\SI{10}{\ohm}$ se
\textbf{abre}. Determine o que muda nas correntes de ramo, na corrente total e
em $R_{eq}$.
\resposta{Ramos restantes inalterados ($\SI{3}{\ampere}$ cada);
$I_t$: $12\to\SI{6}{\ampere}$; $R_{eq}$: $5\to\SI{10}{\ohm}$}
\resolucao{Com fonte ideal, $V=\SI{60}{\volt}$ permanece fixo. Como cada ramo
depende \emph{só} da tensão comum e do próprio resistor, os dois ramos íntegros
continuam com $60/20=\SI{3}{\ampere}$ --- \textbf{nada muda para eles}.\\
O que muda é o agregado: $I_t=3+3=\SI{6}{\ampere}$ (queda de $50\%$) e
$R_{eq}=\SI{10}{\ohm}$.\\
\textbf{Por isso a ligação residencial é em paralelo:} a queima de uma lâmpada
não afeta as demais. Em série, uma falha derruba o circuito inteiro.\\
\emph{Ressalva importante:} isso depende de a fonte ser ideal. Com resistência
interna, $V$ sobe quando um ramo abre --- caso tratado no bloco de desafio.}
\end{questao}

\begin{questao}[3]{}
Um dos ramos de uma associação paralela sofre \textbf{curto-circuito}. Analise
$R_{eq}$, a corrente total e o efeito sobre os demais ramos, e justifique a
necessidade de proteção.
\resposta{$R_{eq}\to0$; $I_t\to\infty$ (limitada só pela fonte e pela fiação);
os demais ramos ficam sem tensão}
\resolucao{Um ramo em curto tem $R=0$, portanto $G\to\infty$. Como
$G_{eq}=\sum G_k$, a condutância total explode e $R_{eq}\to0$:
\[
R_{eq}=\frac{1}{\infty+G_2+G_3}\to0.
\]
Com fonte ideal, $I_t=V/R_{eq}\to\infty$. Na prática, a corrente é limitada
pela resistência interna da fonte e pela impedância da fiação --- mas ainda
assim atinge dezenas ou centenas de ampères.\\
\textbf{Efeito sobre os demais:} o curto derruba a tensão do barramento a
praticamente zero. Os outros ramos, embora íntegros, \emph{param de funcionar}
--- eles são vítimas, não causa.\\
\textbf{Proteção:} é este o cenário que o disjuntor combate. Diferentemente do
ramo aberto (falha benigna e localizada), o curto em um ramo compromete o
sistema inteiro e libera energia suficiente para incendiar a instalação. Daí a
assimetria de tratamento: monitoram-se \emph{sobrecorrentes}, não
subcorrentes.}
\end{questao}

\begin{questao}[3]{}
Dispõe-se de resistores de $\SI{220}{\ohm}$ com dissipação máxima de
$\SI{1}{\watt}$. Projete uma associação paralela de $\SI{55}{\ohm}$ capaz de
dissipar pelo menos $\SI{4}{\watt}$, e determine a tensão máxima aplicável.
\resposta{Quatro em paralelo; $P_{\max}=\SI{4}{\watt}$;
$V_{\max}=\SI{14.8}{\volt}$}
\resolucao{\textbf{Valor:} $220/n=55\Rightarrow n=4$.\\
\textbf{Potência:} sendo iguais e sob a mesma tensão, os quatro dividem a carga
térmica igualmente, então $P_{tot}=4\times1=\SI{4}{\watt}$ \checkmark.\\
\textbf{Tensão máxima:} limita-se pelo componente individual:
\[
V_{\max}=\sqrt{P_1R_1}=\sqrt{1\cdot220}=\SI{14.8}{\volt}.
\]
\textbf{Conferência:} $P_{tot}=V^{2}/R_{eq}=220/55=\SI{4}{\watt}$ \checkmark.\\
\textbf{Observação de projeto:} o paralelo multiplica a capacidade térmica
\emph{e} divide o valor por $n$; a série multiplica ambos. Se o objetivo exige
valor \emph{maior} com mais potência, use série; valor \emph{menor}, paralelo.
Para manter o valor e aumentar a potência, é preciso uma matriz
série--paralelo de $n\times n$ unidades.}
\end{questao}

\begin{questao}[3]{}
Uma rede tem cinco ramos em paralelo: $\SI{10}{\ohm}$, $\SI{20}{\ohm}$,
$\SI{20}{\ohm}$, $\SI{40}{\ohm}$ e $\SI{40}{\ohm}$. Calcule $R_{eq}$ e explique
por que o cálculo em condutâncias é preferível.
\resposta{$G=\SI{0.25}{\siemens}$; $R_{eq}=\SI{4}{\ohm}$}
\resolucao{
\[
G=0{,}100+0{,}050+0{,}050+0{,}025+0{,}025=\SI{0.25}{\siemens}
\;\Longrightarrow\;
R_{eq}=\SI{4}{\ohm}.
\]
\textbf{Três vantagens da condutância.} \emph{(i)} A operação é uma soma
simples, com um único inverso no fim --- contra quatro aplicações sucessivas da
regra produto/soma. \emph{(ii)} A distribuição de corrente sai de graça:
$I_k/I_t=G_k/G_{eq}$, aqui $40\%$, $20\%$, $20\%$, $10\%$ e $10\%$.
\emph{(iii)} Para incerteza, a propagação em $G$ é \emph{linear}, ao passo que
em $R$ envolve derivadas de uma função racional.\\
É também a razão pela qual a análise nodal --- construída sobre a LKC --- adota
condutâncias como grandeza natural.}
\end{questao}

\begin{questao}[3]{}
Um galvanômetro de $\SI{99}{\ohm}$ e $\SI{1}{\milli\ampere}$ de fundo de escala
deve medir até $\SI{100}{\milli\ampere}$.
\begin{itensq}
\item Determine o shunt.
\item Calcule a resistência do amperímetro resultante.
\item Explique por que essa resistência precisa ser pequena.
\end{itensq}
\resposta{(a) $R_{sh}=\SI{1}{\ohm}$; (b) $\SI{0.99}{\ohm}$}
\resolucao{(a) O shunt conduz $100-1=\SI{99}{\milli\ampere}$:
\[
R_{sh}=\frac{R_gI_g}{I_{sh}}=\frac{99\cdot1}{99}=\SI{1}{\ohm}.
\]
(b) $R_A=R_g\parallel R_{sh}=\dfrac{99\cdot1}{100}=\SI{0.99}{\ohm}$.\\
(c) O amperímetro é inserido \textbf{em série} no ramo sob medição, somando-se
à resistência do circuito. Se $R_A$ não for desprezível diante do circuito, a
própria medição altera a corrente que se pretende medir --- o chamado
\emph{erro de inserção}.\\
\textbf{Exemplo:} medindo um ramo de $\SI{100}{\ohm}$, o erro é
$0{,}99/100\approx1\%$; medindo um de $\SI{1}{\ohm}$, seria de $50\%$.\\
Note a dualidade: o \emph{amperímetro} deve ter resistência quase nula e vai em
série; o \emph{voltímetro}, resistência quase infinita e vai em paralelo.}
\end{questao}

\begin{questao}[3]{}
Uma carga de $\SI{100}{\ohm}$ opera a $\SI{200}{\volt}$. A isolação
degradada cria um caminho de fuga de $\SI{10}{\kilo\ohm}$ em paralelo.
\begin{itensq}
\item Determine a corrente de fuga e a corrente total.
\item Calcule $R_{eq}$ e o erro relativo cometido ao ignorar a fuga.
\item Comente o risco.
\end{itensq}
\resposta{(a) $\SI{20}{\milli\ampere}$ e $\SI{2.02}{\ampere}$;
(b) $\SI{99.01}{\ohm}$, erro de $1\%$; (c) risco de choque}
\resolucao{(a) $I_{carga}=200/100=\SI{2}{\ampere}$;
$I_{fuga}=200/10000=\SI{20}{\milli\ampere}$; total $\SI{2.02}{\ampere}$.\\
(b) $R_{eq}=\dfrac{100\cdot10000}{10100}=\SI{99.01}{\ohm}$ --- desvio de
$\approx1\%$ em relação a $\SI{100}{\ohm}$.\\
(c) Do ponto de vista \emph{elétrico}, $1\%$ é irrelevante: nenhum instrumento
de painel acusaria a anomalia. Do ponto de vista da \textbf{segurança}, é
gravíssimo: $\SI{20}{\milli\ampere}$ atravessando um corpo humano já provoca
tetanização muscular, e a partir de $\SI{30}{\milli\ampere}$ o risco de
fibrilação é real.\\
\textbf{Por isso existe o DR (dispositivo diferencial residual):} ele não mede
a corrente total, e sim o \emph{desequilíbrio} entre fase e neutro --- que é
exatamente a corrente de fuga. Um DR de $\SI{30}{\milli\ampere}$ detecta o que
os $\SI{2}{\ampere}$ da carga mascaram completamente.}
\end{questao}

\begin{questao}[3]{}
$R_1=\SI{100}{\ohm}\pm5\%$ e $R_2=\SI{400}{\ohm}\pm5\%$ estão em paralelo.
Determine $R_{eq}$ nominal e seus extremos de pior caso.
\resposta{$R_{eq}=\SI{80}{\ohm}$, entre $\SI{76}{\ohm}$ e $\SI{84}{\ohm}$
($\pm5\%$)}
\resolucao{Nominal: $R_{eq}=\dfrac{100\cdot400}{500}=\SI{80}{\ohm}$.\\
Como $R_{eq}$ cresce monotonicamente com $R_1$ e com $R_2$, os extremos ocorrem
com \emph{ambos} nos extremos:
\[
R_{eq}^{\min}=\frac{95\cdot380}{475}=\SI{76}{\ohm},
\qquad
R_{eq}^{\max}=\frac{105\cdot420}{525}=\SI{84}{\ohm}.
\]
Ou seja, $80\pm4$, isto é, $\pm5\%$ --- exatamente a tolerância das parcelas.\\
\textbf{Resultado notável:} o paralelo \emph{também} preserva a tolerância
relativa, assim como a série. Isso não é coincidência, e a razão é demonstrada
no bloco de desafio.}
\end{questao}

\begin{questao}[3]{}
$R_1=\SI{100}{\ohm}$ com $\alpha=\SI{0.004}{\per\celsius}$ está em paralelo com
$R_2=\SI{100}{\ohm}$ de coeficiente desprezível. Determine $R_{eq}$ a
$\SI{20}{\celsius}$ e a $\SI{80}{\celsius}$, e calcule a variação percentual.
\resposta{$\SI{50}{\ohm}$ e $\SI{55.36}{\ohm}$; variação de $+10{,}7\%$}
\resolucao{A $\SI{20}{\celsius}$: $R_{eq}=100/2=\SI{50}{\ohm}$.\\
A $\SI{80}{\celsius}$: $R_1=100[1+0{,}004(60)]=\SI{124}{\ohm}$ e
\[
R_{eq}=\frac{124\cdot100}{224}=\SI{55.36}{\ohm}.
\]
Variação: $(55{,}36-50)/50=+10{,}7\%$.\\
\textbf{Comparação instrutiva:} $R_1$ variou $24\%$, mas $R_{eq}$ variou
$10{,}7\%$ --- o ramo estável "amortece" a deriva. Curiosamente, o valor é o
mesmo obtido para a \emph{corrente} da associação série equivalente: em ambos
os casos o elemento estável divide o efeito aproximadamente pela metade, porque
as duas resistências são comparáveis.}
\end{questao}

\begin{questao}[3]{}
Um resistor de $\SI{1000}{\ohm}$ está em paralelo com um de $\SI{10}{\ohm}$.
\begin{itensq}
\item Calcule $R_{eq}$.
\item Recalcule supondo que $R_1$ varie $+10\%$.
\item Interprete o resultado em termos de sensibilidade.
\end{itensq}
\resposta{(a) $\SI{9.901}{\ohm}$; (b) $\SI{9.910}{\ohm}$; (c) variação de
apenas $0{,}09\%$}
\resolucao{(a) $R_{eq}=\dfrac{1000\cdot10}{1010}=\SI{9.901}{\ohm}$.\\
(b) Com $R_1=\SI{1100}{\ohm}$:
$R_{eq}=\dfrac{1100\cdot10}{1110}=\SI{9.910}{\ohm}$.\\
(c) Uma variação de $10\%$ em $R_1$ produziu $0{,}09\%$ em $R_{eq}$ ---
atenuação de mais de cem vezes. O ramo dominante ($\SI{10}{\ohm}$) fixa
praticamente sozinho o resultado.\\
\textbf{Uso prático:} é assim que se faz um ajuste \emph{fino}. Colocando um
resistor grande em paralelo com um pequeno, obtém-se um deslocamento pequeno e
controlado --- técnica usada em calibração de shunts e de divisores de
precisão. A quantificação exata dessa sensibilidade aparece no bloco de
desafio.}
\end{questao}

\begin{questao}[3]{}
As cargas da questão residencial ($\SI{27.5}{\ampere}$ totais) são alimentadas
por um circuito cuja fiação tem $\SI{0.2}{\ohm}$ de ida e volta, a partir de um
quadro de $\SI{220}{\volt}$.
\begin{itensq}
\item Determine a queda de tensão na fiação e a tensão nas cargas.
\item Calcule a perda de potência na fiação.
\item Avalie se a queda atende ao limite usual de $4\%$.
\end{itensq}
\resposta{(a) $\SI{5.5}{\volt}$; $\SI{214.5}{\volt}$;
(b) $\SI{151.25}{\watt}$; (c) $2{,}5\%$ --- atende}
\resolucao{(a) $\Delta V=R_fI=0{,}2\cdot27{,}5=\SI{5.5}{\volt}$, logo as cargas
recebem $220-5{,}5=\SI{214.5}{\volt}$.\\
(b) $P_f=R_fI^{2}=0{,}2(27{,}5)^{2}=\SI{151.25}{\watt}$ --- dissipados em puro
aquecimento do cabo.\\
(c) $5{,}5/220=2{,}5\%$, abaixo do limite típico de $4\%$ para circuitos
terminais. \checkmark\\
\textbf{Observação sutil:} a fiação está em \emph{série} com o conjunto, mesmo
que as cargas estejam em paralelo entre si. É por isso que ligar mais uma carga
--- que não afeta as demais em uma rede ideal --- na prática \emph{reduz} a
tensão de todas: a corrente total sobe e a queda na fiação também. O efeito é
visível quando as luzes piscam ao ligar um chuveiro.}
\end{questao}

\begin{questao}[3]{}
Deseja-se obter $\SI{75}{\ohm}$ dispondo de um resistor de $\SI{100}{\ohm}$.
Determine o resistor a ser associado em paralelo e discuta a viabilidade.
\resposta{$R=\SI{300}{\ohm}$}
\resolucao{
\[
\frac{1}{75}=\frac{1}{100}+\frac{1}{R}
\;\Longrightarrow\;
\frac{1}{R}=\frac{1}{75}-\frac{1}{100}=\frac{4-3}{300}
\;\Longrightarrow\;
R=\SI{300}{\ohm}.
\]
\textbf{Viabilidade:} $\SI{300}{\ohm}$ existe na série E24, e o alvo
($\SI{75}{\ohm}$) está a $25\%$ do resistor de partida --- folga confortável.
A precisão do resultado depende principalmente do resistor de
$\SI{100}{\ohm}$, que domina o paralelo.\\
\textbf{Limite do método:} só se pode \emph{reduzir} o valor de partida; alvos
acima de $\SI{100}{\ohm}$ exigem série. E, como se verá no desafio, alvos muito
próximos de $\SI{100}{\ohm}$ tornam a solução impraticável.}
\end{questao}

\begin{questao}[3]{}
Compare a mesma dupla $R_1=\SI{30}{\ohm}$ e $R_2=\SI{60}{\ohm}$ ligada em série
e em paralelo, sob $\SI{90}{\volt}$: resistência, corrente total, distribuição
de corrente e de potência.
\resposta{Série: $\SI{90}{\ohm}$, $\SI{1}{\ampere}$, $P=30$ e
$\SI{60}{\watt}$. Paralelo: $\SI{20}{\ohm}$, $\SI{4.5}{\ampere}$,
$P=270$ e $\SI{135}{\watt}$}
\resolucao{\textbf{Série:} $R=\SI{90}{\ohm}$, $I=\SI{1}{\ampere}$ (comum),
$P_1=30$, $P_2=\SI{60}{\watt}$; total $\SI{90}{\watt}$. O \emph{maior}
resistor dissipa mais.\\
\textbf{Paralelo:} $R=\SI{20}{\ohm}$, $I_1=3$, $I_2=\SI{1.5}{\ampere}$,
$I_t=\SI{4.5}{\ampere}$, $P_1=\SI{270}{\watt}$, $P_2=\SI{135}{\watt}$; total
$\SI{405}{\watt}$. O \emph{menor} resistor dissipa mais.\\
\textbf{Razões:} $R_{s}/R_{p}=90/20=4{,}5$, e a potência total é $4{,}5$ vezes
maior no paralelo --- o mesmo fator, pois $P=V^{2}/R$ com $V$ fixo.\\
\textbf{Síntese:} de dois resistores obtêm-se quatro valores
($30$, $60$, $20$ e $\SI{90}{\ohm}$) e duas distribuições opostas de potência.
A ligação é uma variável de projeto tão importante quanto os componentes.}
\end{questao}

\blocodesafio

\begin{questao}[4]{}
\textbf{(Projeto de instrumentação.)} Um galvanômetro tem $R_g=\SI{99}{\ohm}$ e
fundo de escala $I_g=\SI{1}{\milli\ampere}$. Deseja-se um amperímetro de três
escalas: $\SI{10}{\milli\ampere}$, $\SI{100}{\milli\ampere}$ e
$\SI{1}{\ampere}$.
\begin{itensq}
\item Determine os três shunts na configuração de shunts comutados.
\item Aponte o risco dessa configuração e descreva a alternativa (shunt de
      Ayrton).
\item Calcule a resistência do instrumento em cada escala.
\end{itensq}
\resposta{(a) $\SI{11}{\ohm}$, $\SI{1}{\ohm}$ e $\SI{0.0991}{\ohm}$;
(c) $\SI{9.9}{\ohm}$, $\SI{0.99}{\ohm}$ e $\SI{0.099}{\ohm}$}
\resolucao{(a) Com $n=I_{\max}/I_g$, vale
$R_{sh}=\dfrac{R_g}{n-1}$:
\[
n=10\Rightarrow\frac{99}{9}=\SI{11}{\ohm};\quad
n=100\Rightarrow\frac{99}{99}=\SI{1}{\ohm};\quad
n=1000\Rightarrow\frac{99}{999}=\SI{0.0991}{\ohm}.
\]
(b) \textbf{Risco:} durante a comutação, a chave passa por um instante em que
\emph{nenhum} shunt está conectado. Toda a corrente de linha --- até
$\SI{1}{\ampere}$ --- atravessaria o galvanômetro, projetado para
$\SI{1}{\milli\ampere}$: destruição imediata.\\
\textbf{Shunt de Ayrton:} os resistores ficam permanentemente ligados em série
entre si e o conjunto em paralelo com o galvanômetro; a chave apenas escolhe
\emph{em que ponto} da cadeia a corrente entra. Assim o galvanômetro nunca
fica desprotegido, mesmo durante a comutação.\\
(c) $R_A=R_g\parallel R_{sh}$: $99\parallel11=\SI{9.9}{\ohm}$;
$99\parallel1=\SI{0.99}{\ohm}$; $99\parallel0{,}0991=\SI{0.099}{\ohm}$.\\
Note que $R_A=R_g/n$ exatamente --- escalas maiores exigem instrumentos de
resistência menor, o que é justamente o desejável.}
\end{questao}

\begin{questao}[4]{}
\textbf{(Demonstração --- escalonamento de bancos.)} Considere $N$ resistores
iguais de valor $R$ e potência nominal $P_1$.
\begin{itensq}
\item Mostre que a associação paralela vale $R/N$ e suporta $NP_1$.
\item Mostre que a série vale $NR$ e também suporta $NP_1$.
\item Deduza a matriz $m\times n$ necessária para manter o valor $R$ e obter
      potência $mnP_1$.
\end{itensq}
\resposta{(c) $m$ colunas em paralelo de $n$ unidades em série, com $m=n$}
\resolucao{(a) $G_{eq}=NG\Rightarrow R_{eq}=R/N$. Sob tensão $V$ comum, cada
unidade dissipa $V^{2}/R$; como todas são iguais e submetidas à mesma tensão, o
total é $N$ vezes o individual: $P_{\max}=NP_1$.\\
(b) $R_{eq}=NR$. Com corrente $I$ comum, cada unidade dissipa $RI^{2}$ e o
total é novamente $NP_1$.\\
\textbf{Fato central:} \emph{qualquer} associação de $N$ resistores iguais
suporta $NP_1$, desde que a dissipação se reparta igualmente. O que a topologia
escolhe é o \emph{valor}, não a capacidade térmica.\\
(c) Uma matriz de $n$ unidades em série, repetida em $m$ colunas paralelas:
\[
R_{eq}=\frac{nR}{m},\qquad P_{\max}=mnP_1.
\]
Para preservar o valor, $n=m$; então $R_{eq}=R$ e
$P_{\max}=n^{2}P_1$.\\
\textbf{Exemplo:} para quadruplicar a potência sem alterar o valor, use
$2\times2=4$ unidades. Para $\times9$, use $3\times3=9$. A potência cresce com
o \emph{quadrado} do lado da matriz --- e o custo, também.}
\end{questao}

\begin{questao}[4]{}
\textbf{(Sensibilidade.)} Para $R_{eq}=\dfrac{R_1R_2}{R_1+R_2}$:
\begin{itensq}
\item Obtenha $\dfrac{\partial R_{eq}}{\partial R_1}$ e mostre que vale
      $\left(\dfrac{R_2}{R_1+R_2}\right)^{2}$.
\item Obtenha a sensibilidade relativa e mostre que
      $S_1+S_2=1$.
\item Avalie para $R_1=\SI{100}{\ohm}$ e $R_2=\SI{400}{\ohm}$, e interprete.
\end{itensq}
\resposta{(a) $0{,}64$; (b) $S_1=R_2/(R_1+R_2)$, $S_2=R_1/(R_1+R_2)$;
(c) $S_1=0{,}8$ e $S_2=0{,}2$}
\resolucao{(a) Derivando o quociente em relação a $R_1$:
\[
\frac{\partial R_{eq}}{\partial R_1}
=\frac{R_2(R_1+R_2)-R_1R_2}{(R_1+R_2)^{2}}
=\frac{R_2^{2}}{(R_1+R_2)^{2}}
=\left(\frac{R_2}{R_1+R_2}\right)^{2}.
\]
Com $R_1=100$ e $R_2=400$: $(400/500)^{2}=0{,}64$ --- adimensional, pois é
$\si{\ohm}$ por $\si{\ohm}$.\\
(b)
\[
S_1=\frac{\partial R_{eq}}{\partial R_1}\cdot\frac{R_1}{R_{eq}}
=\frac{R_2^{2}}{(R_1+R_2)^{2}}\cdot\frac{R_1(R_1+R_2)}{R_1R_2}
=\frac{R_2}{R_1+R_2}.
\]
Por simetria, $S_2=\dfrac{R_1}{R_1+R_2}$, e a soma é $1$ --- consequência de
$R_{eq}$ ser \textbf{homogênea de grau 1}.\\
(c) $S_1=0{,}8$ e $S_2=0{,}2$: um erro de $1\%$ em $R_1$ produz $0{,}8\%$ em
$R_{eq}$, e o mesmo erro em $R_2$ produz apenas $0{,}2\%$.\\
\textbf{Regra de projeto:} o resistor \emph{menor} domina o paralelo e deve
receber a tolerância apertada --- exatamente o inverso da série, onde manda o
maior. E se $R_2\gg R_1$, então $S_1\to1$ e $S_2\to0$: o ramo grande vira
ajuste fino, o que confirma quantitativamente o observado na questão do
$\SI{1000}{\ohm}$ com $\SI{10}{\ohm}$.}
\end{questao}

\begin{questao}[4]{}
\textbf{(Divisão de corrente sob tolerância.)} Três shunts iguais de
$\SI{0.01}{\ohm}$ devem dividir igualmente $\SI{30}{\ampere}$.
\begin{itensq}
\item Determine a corrente ideal por shunt e a tensão sobre eles.
\item Se um deles for $2\%$ menor e os outros dois nominais, calcule o
      desequilíbrio de corrente.
\item Discuta a implicação térmica.
\end{itensq}
\resposta{(a) $\SI{10}{\ampere}$ e $\SI{0.1}{\volt}$;
(b) $\SI{10.13}{\ampere}$ contra $\SI{9.93}{\ampere}$ --- desvio de $1{,}3\%$;
(c) o mais "fraco" aquece mais}
\resolucao{(a) $I=30/3=\SI{10}{\ampere}$;
$V=0{,}01\cdot10=\SI{0.1}{\volt}$.\\
(b) Com $R_a=\SI{0.0098}{\ohm}$ e $R_b=R_c=\SI{0.0100}{\ohm}$:
\[
G=\frac{1}{0{,}0098}+\frac{2}{0{,}0100}=102{,}04+200{,}00=\SI{302.04}{\siemens},
\]
\[
I_a=30\cdot\frac{102{,}04}{302{,}04}=\SI{10.13}{\ampere},
\qquad
I_b=I_c=30\cdot\frac{100{,}00}{302{,}04}=\SI{9.93}{\ampere}.
\]
\textbf{Verificação:} $10{,}13+9{,}93+9{,}93=\SI{29.99}{\ampere}$ \checkmark.\\
\textbf{Fato tranquilizador:} um desvio de $2\%$ na resistência produziu apenas
$1{,}3\%$ de desequilíbrio --- a divisão é \emph{mais estável} que os
componentes, porque os outros dois ramos "puxam" a média.\\
(c) $P_a=0{,}0098(10{,}13)^{2}=\SI{1.006}{\watt}$ contra
$P_b=0{,}0100(9{,}93)^{2}=\SI{0.986}{\watt}$: o shunt de menor resistência
dissipa $2\%$ a mais. Se tiver coeficiente térmico positivo, aquecerá,
aumentará sua resistência e \emph{devolverá} corrente aos vizinhos --- uma
realimentação \textbf{estabilizadora}. Com coeficiente negativo ocorreria o
oposto: fuga térmica. É por isso que resistores para divisão de corrente são
feitos em ligas de $\alpha$ positivo e pequeno, como a manganina.}
\end{questao}

\begin{questao}[4]{}
\textbf{(Compensação térmica em paralelo.)} Deseja-se uma associação paralela
com resistência independente da temperatura, usando
$R_1=\SI{100}{\ohm}$ com $\alpha_1=+4000\,\mathrm{ppm/^{\circ}C}$ e um segundo
resistor com $\alpha_2=-1000\,\mathrm{ppm/^{\circ}C}$.
\begin{itensq}
\item Deduza a condição de compensação em termos de condutâncias.
\item Determine $R_2$ e $R_{eq}$.
\item Verifique a compensação e comente sua exatidão.
\end{itensq}
\resposta{(a) $G_1\alpha_1+G_2\alpha_2=0$; (b) $R_2=\SI{25}{\ohm}$,
$R_{eq}=\SI{20}{\ohm}$; (c) exata só em primeira ordem}
\resolucao{(a) Partindo de $G_{eq}=G_1+G_2$ e de
$R_k(T)=R_k(1+\alpha_k\Delta T)$, tem-se
$G_k\approx G_k(1-\alpha_k\Delta T)$ para $\alpha_k\Delta T\ll1$. Então
\[
G_{eq}(T)\approx(G_1+G_2)-(G_1\alpha_1+G_2\alpha_2)\Delta T,
\]
e a condição de estabilidade é
\[
G_1\alpha_1+G_2\alpha_2=0
\quad\Longleftrightarrow\quad
\frac{\alpha_1}{R_1}+\frac{\alpha_2}{R_2}=0.
\]
(b) $R_2=-\dfrac{\alpha_2R_1}{\alpha_1}=\dfrac{1000\cdot100}{4000}
=\SI{25}{\ohm}$, e $R_{eq}=100\parallel25=\SI{20}{\ohm}$.\\
(c) \emph{Para $\Delta T=\SI{1}{\celsius}$:}
$R_1=\SI{100.4}{\ohm}$, $R_2=\SI{24.975}{\ohm}$, $R_{eq}=\SI{20.000}{\ohm}$ ---
compensação praticamente perfeita.\\
\emph{Para $\Delta T=\SI{50}{\celsius}$:} $R_1=\SI{120}{\ohm}$,
$R_2=\SI{23.75}{\ohm}$, $R_{eq}=\SI{19.83}{\ohm}$ --- desvio de $0{,}9\%$.\\
A compensação é de \textbf{primeira ordem}: excelente perto da temperatura de
projeto, degradando-se quadraticamente conforme se afasta.\\
\textbf{Compare com a série}, cuja condição era $R_1\alpha_1+R_2\alpha_2=0$
(ponderada por resistências). Aqui a ponderação é por \emph{condutâncias} ---
mais uma manifestação da dualidade série/paralelo.}
\end{questao}

\begin{questao}[4]{}
\textbf{(Falha de ramo com fonte real.)} Uma fonte de $\SI{24}{\volt}$ com
$r=\SI{2}{\ohm}$ alimenta três ramos de $\SI{6}{\ohm}$ em paralelo.
\begin{itensq}
\item Determine a tensão do barramento, a corrente total e a de cada ramo.
\item Um ramo se abre. Recalcule tudo.
\item Explique por que os ramos remanescentes \emph{mudam}, ao contrário do
      caso com fonte ideal.
\end{itensq}
\resposta{(a) $\SI{12}{\volt}$, $\SI{6}{\ampere}$, $\SI{2}{\ampere}$ por ramo;
(b) $\SI{14.4}{\volt}$, $\SI{4.8}{\ampere}$, $\SI{2.4}{\ampere}$ por ramo}
\resolucao{(a) $R_{eq}=6/3=\SI{2}{\ohm}$;
\[
I_t=\frac{24}{2+2}=\SI{6}{\ampere},\qquad
V_{barra}=24-2(6)=\SI{12}{\volt},\qquad
I_{ramo}=\frac{12}{6}=\SI{2}{\ampere}.
\]
(b) Com dois ramos, $R_{eq}=\SI{3}{\ohm}$:
\[
I_t=\frac{24}{2+3}=\SI{4.8}{\ampere},\qquad
V_{barra}=24-2(4{,}8)=\SI{14.4}{\volt},\qquad
I_{ramo}=\frac{14{,}4}{6}=\SI{2.4}{\ampere}.
\]
(c) Os ramos remanescentes \emph{não mudaram de valor}, mas a tensão que os
alimenta subiu de $12$ para $\SI{14.4}{\volt}$ ($+20\%$), e com ela a corrente
de cada um ($+20\%$).\\
\textbf{Mecanismo:} $r$ está em série com o conjunto. Menos ramos $\Rightarrow$
menor corrente total $\Rightarrow$ menor queda interna $\Rightarrow$ maior
tensão disponível.\\
\textbf{Implicação de projeto:} o isolamento entre ramos paralelos é uma
\emph{idealização}. Com fonte real, ramos interagem via $r$ --- e a perda de
uma carga pode sobrecarregar as demais em $20\%$ ou mais. Fontes reguladas
existem justamente para levar $r$ efetiva a valores próximos de zero e
restaurar o desacoplamento.}
\end{questao}

\begin{questao}[4]{}
\textbf{(Limite do método de ajuste.)} Dispõe-se de um resistor fixo de
$\SI{100}{\ohm}$ e deseja-se obter valores-alvo por associação paralela.
\begin{itensq}
\item Deduza $R(alvo)$ e calcule para alvos de $\SI{90}{\ohm}$,
      $\SI{95}{\ohm}$ e $\SI{99}{\ohm}$.
\item Obtenha a sensibilidade do alvo em relação a $R$ e avalie no caso de
      $\SI{99}{\ohm}$.
\item Conclua sobre a viabilidade prática.
\end{itensq}
\resposta{(a) $\SI{900}{\ohm}$, $\SI{1900}{\ohm}$ e $\SI{9900}{\ohm}$;
(c) inviável para alvos próximos do resistor fixo}
\resolucao{(a) De $\dfrac{1}{R_{alvo}}=\dfrac{1}{R_f}+\dfrac{1}{R}$:
\[
R=\frac{R_fR_{alvo}}{R_f-R_{alvo}}.
\]
\[
90\to\frac{100\cdot90}{10}=\SI{900}{\ohm};\quad
95\to\frac{100\cdot95}{5}=\SI{1900}{\ohm};\quad
99\to\frac{100\cdot99}{1}=\SI{9900}{\ohm}.
\]
O denominador $R_f-R_{alvo}$ tende a zero e $R$ \textbf{diverge}.\\
(b) Pela sensibilidade relativa do paralelo,
$S=\dfrac{R}{R_f+R}$. Para $R=\SI{9900}{\ohm}$:
$S=9900/10000=0{,}99$.\\
Ou seja: $1\%$ de erro em $R$ produz $0{,}99\%$ de erro no alvo. Como o ajuste
pretendido é de apenas $1\%$ (de $100$ para $\SI{99}{\ohm}$), \emph{o erro tem
a mesma ordem de grandeza do ajuste} --- ele o anula.\\
(c) \textbf{Inviável.} Com resistor de $5\%$ de tolerância, o "ajuste" de
$\SI{1}{\ohm}$ vem com incerteza de $\SI{0.99}{\ohm}$: o alvo pode ficar em
qualquer ponto entre $98$ e $\SI{100}{\ohm}$. Além disso,
$\SI{9900}{\ohm}$ não é valor comercial, e a resistência de contato já é
comparável ao efeito buscado.\\
\textbf{Regra prática:} ajuste por paralelo só é confiável para correções de
\emph{dezenas} de por cento. Para correções finas, use potenciômetro de ajuste
ou selecione o componente por medição direta.}
\end{questao}

\begin{questao}[4]{}
\textbf{(Demonstração --- propagação de incerteza.)} Mostre que, se
\emph{todos} os resistores de uma rede puramente resistiva tiverem a mesma
tolerância relativa $t$, então $R_{eq}$ terá exatamente a tolerância $t$,
qualquer que seja a topologia.
\begin{itensq}
\item Demonstre a propriedade de homogeneidade.
\item Verifique em um exemplo série e em um paralelo.
\item Discuta o que muda quando as tolerâncias são diferentes.
\end{itensq}
\resposta{$R_{eq}$ é homogênea de grau 1: $R_{eq}(\lambda R_k)=\lambda
R_{eq}(R_k)$; logo a tolerância relativa se conserva}
\resolucao{(a) Qualquer $R_{eq}$ de rede resistiva é construída por somas
($R_i+R_j$) e por combinações do tipo $\dfrac{R_iR_j}{R_i+R_j}$ --- ambas
\textbf{homogêneas de grau 1}. Multiplicando \emph{todos} os resistores por
$\lambda$:
\[
\lambda R_i+\lambda R_j=\lambda(R_i+R_j),
\qquad
\frac{(\lambda R_i)(\lambda R_j)}{\lambda R_i+\lambda R_j}
=\lambda\frac{R_iR_j}{R_i+R_j}.
\]
Por indução sobre a construção da rede, $R_{eq}(\lambda R_k)=\lambda R_{eq}$.
Fazendo $\lambda=1\pm t$, obtém-se $R_{eq}(1\pm t)$: tolerância relativa
exatamente $t$.\\
(b) \emph{Série:} $100\pm5\%$ e $200\pm5\%$ dão $300\pm\SI{15}{\ohm}=\pm5\%$
\checkmark. \emph{Paralelo:} $100\pm5\%$ e $400\pm5\%$ dão
$80\pm\SI{4}{\ohm}=\pm5\%$ \checkmark. A demonstração explica por que os dois
casos, aparentemente distintos, deram o mesmo resultado.\\
(c) Com tolerâncias \emph{diferentes}, a homogeneidade não se aplica --- os
resistores não escalam pelo mesmo $\lambda$. Vale então a combinação ponderada
pelas sensibilidades:
\[
t_{eq}=\sum_k S_k\,t_k,
\qquad \text{com } \sum_k S_k=1.
\]
Em série, $S_k=R_k/\sum R_j$ (manda o maior); em paralelo,
$S_k=G_k/\sum G_j$ (manda o menor). A restrição $\sum S_k=1$ é justamente a
identidade de Euler para funções homogêneas de grau 1 --- e garante que
$t_{eq}$ jamais ultrapasse a maior das tolerâncias individuais.}
\end{questao}

\begin{questao}[4]{}
\textbf{(Seletividade da proteção.)} Um barramento de $\SI{220}{\volt}$
alimenta três circuitos em paralelo, com correntes nominais de
$\SI{20}{\ampere}$, $\SI{5}{\ampere}$ e $\SI{2.5}{\ampere}$, protegidos
individualmente e por um disjuntor geral.
\begin{itensq}
\item Especifique os disjuntores de ramo e o geral.
\item Explique o requisito de seletividade e o que ocorre se ele falhar.
\item Discuta por que a soma dos disjuntores de ramo pode exceder o geral.
\end{itensq}
\resposta{(a) ramos de $25$, $10$ e $\SI{10}{\ampere}$; geral de
$\SI{32}{\ampere}$; (b) o ramo deve atuar antes do geral}
\resolucao{(a) Cada disjuntor de ramo fica acima da corrente nominal do próprio
circuito: $\SI{25}{\ampere}$, $\SI{10}{\ampere}$ e $\SI{10}{\ampere}$
(valores comerciais). O geral cobre a soma das cargas,
$\SI{27.5}{\ampere}$: disjuntor de $\SI{32}{\ampere}$.\\
(b) \textbf{Seletividade} significa que, diante de uma falta em um ramo,
\emph{apenas} o disjuntor daquele ramo abre --- os demais circuitos continuam
alimentados. Isso exige que a curva tempo--corrente do dispositivo de ramo
esteja inteiramente \emph{abaixo} da do geral.\\
Se a seletividade falhar, um curto em uma tomada derruba a instalação inteira:
além do transtorno, perde-se a informação de \emph{onde} está o defeito, e a
localização passa a exigir religamentos sucessivos.\\
(c) Porque as cargas dificilmente operam \emph{todas} no máximo ao mesmo tempo
--- é o \textbf{fator de demanda}. Aqui $25+10+10=\SI{45}{\ampere}$ contra um
geral de $\SI{32}{\ampere}$: perfeitamente admissível, desde que a demanda
simultânea real permaneça abaixo de $\SI{32}{\ampere}$.\\
\textbf{Consequência do paralelo:} como as correntes de ramo se somam no
alimentador, é o \emph{comportamento estatístico} do conjunto --- e não a soma
aritmética dos máximos --- que dimensiona a entrada. Superdimensionar o geral
para $\SI{45}{\ampere}$ exigiria condutores muito mais caros sem ganho real de
disponibilidade.}
\end{questao}
